\documentclass[12pt]{article}
\usepackage[utf8]{inputenc}
\usepackage[T1]{fontenc}
\usepackage{hyperref}
\title{Ho-Chunk Language Resources for K-12 Curriculum Development}
\author{Compiled by Claude Sonnet 5.5 (claude-sonnet-5-5) for Andr\'as Kornai}
\date{}
\begin{document}
\maketitle

\section{Introduction}
Ho-Chunk (Hoc\k{a}k; ISO 639-3: win; Glottocode hoch1243) is a Siouan language of the United States. The Ho-Chunk Nation is based in Wisconsin, and Black River Falls is mentioned in the search results. This report is based on a single web search. Only resources that appeared in the search results are listed.

\section{Primary Educational Resources}
The search found no specific K-12 curriculum document. It did find learner-oriented materials that could serve as a basis for one:
\begin{itemize}
\item \emph{Hoc\k{a}k Teaching Materials, Volume 1: Elements of Grammar / Learner's Dictionary} (SUNY Press, 2010; earlier edition 2006), edited by Johannes Helmbrecht and Christian Lehmann with the Hoc\k{a}k Wazija Haci Language Division. \url{https://wiki.ubc.ca/Documentation:RelLex/Ho%C3%A7ak_Teaching_Materials,_Volume_1:_Elements_of_Grammar/Learner%27s_Dictionary}
\item \emph{Hocak Teaching Materials Volume 12}, a SUNY Press publication. \url{https://sunypress.edu/Books/H/Hocak-Teaching-Materials-Volume-12}
\item \emph{Ho-Chunk Phrasebook} (2024) by Mateo Russo and Sandra Chiguela. It includes a grammar section, an English--Ho-Chunk dictionary and everyday phrases. \url{https://shop.wisconsinhistory.org/ho-chunk-phrasebook}
\end{itemize}

\section{Linguistic and Reference Materials}
\begin{itemize}
\item \emph{Wisconsin Winnebago Lexicon} (1976), a bilingual, bidirectional word list. \url{https://wiki.ubc.ca/Documentation:RelLex/Wisconsin_Winnebago_Lexicon}
\item \emph{A Lexical Study of Winnebago} (1988), based on Josephine P. White Eagle's PhD thesis. \url{https://wiki.ubc.ca/Documentation:RelLex/A_Lexical_Study_of_Winnebago}
\item Library catalogue record for a related volume. \url{https://library.holycross.edu/Record/eres_ocn870962550}
\item A scholarly PDF on Ho-Chunk. Its title and content were not described in the results. \url{https://www.christianlehmann.eu/publ/ASSidUE21.pdf}
\end{itemize}

\section{Religious and Community Resources}
No scripture or community-organisation URLs appeared in the results.

\section{Digital and Audio Resources}
The Ho-Chunk Dictionary mobile app was built by The Language Conservancy. It works offline, has audio, and lists 11,961 entries. It was developed with more than 20 native speakers. \url{https://apps.apple.com/app/id1633317298}

\section{Curriculum Development Context}
Ho-Chunk is endangered and spoken mainly by older adults. It is not known to be taught in schools. The Teaching Materials series and the mobile dictionary with audio are the strongest starting points. A full curriculum would need collaboration with the Ho-Chunk Nation's language programme. This search did not find a programme page, so one should be located separately.

\end{document}